\documentclass[10pt,a4paper]{article}

\usepackage[utf8]{inputenc}
\usepackage[T1]{fontenc}
\usepackage[margin=2.5cm]{geometry}
\usepackage{setspace}
\usepackage{parskip}
\usepackage{enumitem}
\usepackage{hyperref}
\usepackage{xcolor}
\usepackage{tabularx}
\usepackage{array}
\usepackage{seqsplit}
\usepackage{footmisc}

\hypersetup{
    colorlinks=true,
    linkcolor=black,
    urlcolor=blue,
    citecolor=black
}

\setstretch{1.00}
\setlist[itemize]{topsep=4pt,itemsep=2pt,leftmargin=1.8em}

\title{\textbf{Generating Wikipedia Articles from Citations}}
\author{Nino Reil, Andreas Wallnöfer, Zoltan Bek, Armin Lohse}
\date{\today}

\begin{document}

\maketitle

\section{Problem Description}

Wikipedia articles are usually written as concise, neutral, and well-structured summaries of knowledge that is grounded in external sources. In practice, the references of a Wikipedia article often point to books, news articles, reports, scientific papers, or official websites that support individual claims in the text. This project asks the question: \emph{Can we reconstruct a plausible Wikipedia-style article if we are only given the cited source material, but not the original article itself?}

A central challenge is that the cited sources may differ strongly in tone, structure, and writing style, while the generated output should follow the neutral, professional, and objective style typical of Wikipedia articles. The system must therefore not only extract relevant information from the sources, but also transform it into a consistent encyclopedic form.

\section{Dataset Description}

We will build the dataset from existing Wikipedia articles and their references.

\begin{itemize}
    \item \textbf{Dataset Source:} Wikipedia articles together with their cited external references.
    \item \textbf{Type of Text:} Encyclopedic article text, citation metadata, and source documents such as news articles, reports, official websites, and scientific abstracts.
    \item \textbf{Approximate Size:} About 420 Wikipedia articles.
    \item \textbf{Language:} English.
    \item \textbf{Public or Collected:} Collected from public sources.
\end{itemize}

To keep the task feasible and the source coverage reliable, we will only include Wikipedia articles that satisfy two conditions:
\begin{itemize}
    \item At least 20 citations in total.
    \item At least 90\% of these citations are accessible to us, for example because they are not paywalled or missing.
\end{itemize}

Each sample consists of one Wikipedia article and its citation set. The original article text is used only as a reference target for evaluation, not as input to the generation system.

This is mainly a generation task, so the dataset does not contain class labels in the standard classification sense. Instead, each sample has:
\begin{itemize}
    \item \textbf{Input:} Cited source documents.
    \item \textbf{Target:} The original Wikipedia article text.
\end{itemize}

\section{Methodology Approach}

The dataset of 420 samples will be divided into
\begin{enumerate}
    \item \textbf{Training:} 300
    \item \textbf{Validation:} 60
    \item \textbf{Test:} 60
\end{enumerate}

The validation set will be used for model selection and parameter tuning, while the test set will be kept untouched for the final evaluation.

\subsection*{Shared Pipeline}

Both approaches use the same initial processing steps:

\begin{enumerate}
    \item \textbf{Citation Extraction and Document Collection}

    Citations are extracted from selected Wikipedia articles, and accessible referenced documents are downloaded automatically.

    \item \textbf{Document Cleaning and Passage Creation}

    Web noise such as navigation text, advertisements, and unrelated page elements is removed. The main textual content is extracted and split into passages of approximately 150--300 words.

    \item \textbf{Topic Inference}

    The generation system does not receive the original Wikipedia article title as input. Instead, the topic is inferred automatically from the cited documents. Candidate topic phrases are obtained from document titles, named entities, and frequently occurring noun phrases. These candidates are ranked using frequency and semantic similarity, and the most central candidate is selected as the inferred topic.

    \item \textbf{Passage Embedding and Indexing}

    Each passage is converted into a vector representation using the embedding model
    \texttt{\seqsplit{sentence-transformers/all-MiniLM-L6-v2}}\footnote{All suggested models used in this project were chosen because they are recent and can be locally run with modest hardware requirements.}. The resulting embeddings are stored in a vector index (e.g.\ FAISS) for efficient similarity search.
\end{enumerate}

After these shared steps, the two approaches differ in how they retrieve evidence and generate the final article.

\subsection*{Approach Comparison}

\renewcommand{\arraystretch}{1.2}
\begin{tabularx}{\textwidth}{|>{\raggedright\arraybackslash}p{0.18\textwidth}|>{\raggedright\arraybackslash}X|>{\raggedright\arraybackslash}X|}
\hline
\textbf{} & \textbf{Approach 1: Baseline} & \textbf{Approach 2: Retrieval-Augmented Generation} \\
\hline
\textbf{Retrieval} &
Using the inferred topic, the system retrieves the top-$k$ most relevant passages from the cited documents using cosine similarity in the embedding space. &
The inferred topic is expanded into related subtopics (e.g.\ via noun phrase extraction or keyword expansion). Relevant passages are retrieved separately for each subtopic. Highly similar passages may be filtered using cosine similarity or clustering to reduce redundancy. \\
\hline
\textbf{Generation} &
The retrieved passages serve as evidence for article generation. A compact language model synthesizes the information into a Wikipedia-style article, combining facts from multiple passages and rewriting them in a neutral encyclopedic style. &
Instead of generating the article in one step, passages are grouped by subtopic and used to generate sections separately. For each group, the model writes a coherent Wikipedia-style section grounded in the retrieved evidence. \\
\hline
\textbf{Output Assembly} &
The generated text is treated as the reconstructed article. &
The generated sections are combined into a final reconstructed article. \\
\hline
\end{tabularx}

\vspace{0.5em}

Suitable local models for the generation step in both approaches are:
\begin{itemize}
    \item \texttt{google/gemma-2-2b}\footnote{\label{ftnt:models}All suggested models used in this project were chosen because they are recent and can be locally run with modest hardware requirements.}
    \item \texttt{microsoft/Phi-3-mini-4k-instruct}\footref{ftnt:models}
\end{itemize}

\section{Evaluation Metrics}

The models will be evaluated automatically against the original Wikipedia article, which is not provided as input to the system.

Since exact wording may differ even in a good reconstruction, we will not rely on a single metric. We plan to compare the generated article with the original Wikipedia article using

\begin{enumerate}
    \item \textbf{Title Similarity:} Cosine similarity between embeddings.
    \item \textbf{Lexical Content Similarity:} ROUGE-1, ROUGE-2, and ROUGE-L.
    \item \textbf{Semantic Content Similarity:} BERTScore.
    \item \textbf{Absolute Difference in Section Count}
    \item \textbf{Article Length Ratio}
\end{enumerate}

\section{Expected Results}

We expect the retrieval-augmented generation pipeline to outperform the baseline because it can retrieve more relevant information from the cited documents and generate more coherent text.

One of the main challenges will be bridging the gap between heterogeneous source language and the neutral, professional, and objective wording expected in Wikipedia articles. We expect the generated articles to recover the main facts and overall topic, although wording, structure, and level of detail may differ from the original Wikipedia article.

\end{document}